\begin{abstract}

Ẩm thực Việt Nam đa dạng theo vùng miền, trong khi người dùng đôi khi chỉ có
ảnh món ăn mà không biết tên để tìm công thức phù hợp. Đồ án xây dựng hệ
thống Vị Việt nhằm kết hợp khám phá công thức có nguồn với nhận diện món ăn
từ ảnh trên thiết bị di động. Ứng dụng Flutter giao tiếp với backend FastAPI
qua HTTP; dữ liệu nghiệp vụ được lưu bằng SQLite và SQLAlchemy. Ảnh JPEG hoặc
PNG được xử lý bởi mô hình Vision Transformer ViT-L/16 tiền huấn luyện, sau
đó nhãn dự đoán được chuẩn hóa và ánh xạ tới bản ghi món ăn trong cơ sở dữ
liệu. Hệ thống hiện quản lý 38 món đã xuất bản, 38 ảnh có metadata giấy phép,
348 nguyên liệu, 205 bước nấu và 85 nguồn tham khảo. Tám công thức được dịch
và chuẩn hóa từ nguồn mở; ba mươi công thức còn lại do Vị Việt biên soạn và
đối chiếu tối thiểu hai nguồn. Bộ dữ liệu 30VNFoods gồm 25.136 ảnh thuộc 30
lớp được audit và tạo manifest sạch còn 25.032 ảnh; tập kiểm thử sạch gồm
5.014 ảnh được dùng để đánh giá độc lập 30 lớp giao nhau với model. Ngoài
nhận diện, ứng dụng hỗ trợ tìm kiếm, lọc, yêu thích, lịch sử, phản hồi kết
quả và hiển thị nguồn--bản quyền. Hệ thống chứng minh khả năng kết hợp mô
hình thị giác tiền huấn luyện với một kho công thức có thể truy nguyên trong
một ứng dụng mobile hoàn chỉnh, đồng thời chỉ rõ giới hạn khi tập đánh giá
chưa bao phủ đủ 38 nhãn của model.

\vspace{0.8em}
\noindent\textbf{Từ khoá:} nhận diện món ăn, Vision Transformer, Flutter,
FastAPI, công thức món Việt, 30VNFoods.

\end{abstract}
